% \iffalse meta-comment
%
% Copyright (C) 2017- by Yanshuo Chu <yanshuoc@gmail.com>
%
% This file may be distributed and/or modified under the
% conditions of the LaTeX Project Public License, either version 1.3a
% of this license or (at your option) any later version.
% The latest version of this license is in:
%
% http://www.latex-project.org/lppl.txt
%
% and version 1.3a or later is part of all distributions of LaTeX
% version 2004/10/01 or later.
%
% \fi
%
% \section{正文的字体与字号}
% \changes{v3.2a}{2026/09/13}{终稿与报告的字号并进一份，主体在 \file{src/engine/hit-fonts.dtx}；
%   正文字号与行距从 \cs{g\_docxlike\_body\_size\_dim}、\cs{g\_docxlike\_body\_lead\_dim} 取，
%   这里只给默认值}
%
% ^^A ======================================================================
% ^^A Module flow-fonts: 正文字号、中英文主字体与字号宏命令（\wuhao 等）
% ^^A ======================================================================
%    \begin{macrocode}
%<*flow-fonts>
%    \end{macrocode}
%
% 正文的字号与行距。两个数在这里只是默认值，真正的取值由 \option{styles} 的
% \option{Normal} 样式算出来，机制见 \file{src/docxlike/docxlike-format.dtx} 的「接到正文上」。
% \cs{normalsize} 每次执行都重读这两个长度，所以导言区里改得动。
%
% 字号两档一样，行距与 \cs{normalsize} 分档，所以下面按 \option{stage} 分两段。
%    \begin{macrocode}
\hit@setup@fonts
\dim_gset:Nn \g_docxlike_body_size_dim { 12bp }
%    \end{macrocode}
%
% \subsection{终稿的行距}
%
% 20.50394\,bp 这个默认行距是旧版面的，来路是 v3.1e 的写法，与规范里
% 「多倍行距 1.25」算出来的 19.455\,bp 不是一个数。新版面下会被盖掉。
%
% \cs{hit@body@lead@skip:} 只有终稿这一支调。它读
% \cs{g\_\_hit\_layout\_two\_bool}，而报告没有 \option{v2-layout} 这个选项，
% 那个标志位在报告里不存在，所以报告那一段自己写伸缩量，不经过它。
%    \begin{macrocode}
\bool_if:NT \g__hit_final_bool
  {
    \dim_gset:Nn \g_docxlike_body_lead_dim { 20.50394bp }
    \cs_new:Npn \hit@body@lead@skip:
      {
        \bool_if:NTF \g__hit_layout_two_bool
          { \hit@glue@skip{\dim_use:N \g_docxlike_body_lead_dim}{2.834646bp} }
          {
            \dim_use:N \g_docxlike_body_lead_dim
            \dim_compare:nNnF { \g_docxlike_body_stretch_dim } = { 0pt }
              { ~ plus ~ \dim_use:N \g_docxlike_body_stretch_dim }
          }
      }
    \cs_set:Npn \normalsize
      {
        \@setfontsize \normalsize { \g_docxlike_body_size_dim }
          { \hit@body@lead@skip: }
        \hit@display@skips@auto:n { 8pt }
      }
  }
%    \end{macrocode}
%
% \subsection{开题与中期报告的行距}
%
% 深圳博士中期那一份行距是 15\,bp 且不给伸缩量，其余 20.50394\,bp。
%    \begin{macrocode}
\bool_if:NF \g__hit_final_bool
  {
    \hit@if@shenzhen@doctor@interim@TF
      {
        \dim_gset:Nn \g_docxlike_body_lead_dim { 15bp }
        \cs_set:Npn \normalsize
          {
            \@setfontsize \normalsize { \g_docxlike_body_size_dim }
              { \g_docxlike_body_lead_dim }
            \hit@display@skips@auto:n { 6pt }
          }
      }
      {
        \dim_gset:Nn \g_docxlike_body_lead_dim { 20.50394bp }
        \cs_set:Npn \normalsize
          {
            \@setfontsize \normalsize { \g_docxlike_body_size_dim }
              { \hit@glue@skip{\dim_use:N \g_docxlike_body_lead_dim}{2.834646bp} }
            \hit@display@skips@auto:n { 8pt }
          }
      }
  }
%</flow-fonts>
%    \end{macrocode}
%
% \Finale
